\section{Worked example: proxy ARP}\label{sec:example}
This section follows one ping through the \code{proxy\_arp} preset, event by event. Host h1 pings h3 (141.23.56.21) while every ARP cache is empty. The trace comes from a real run of the system, so every frame matches events in the log.

% Styles for this section only (prefix ex).
\tikzset{
  exhead/.style={draw=black!70, rounded corners=3pt, fill=white, font=\footnotesize,
                 minimum width=1.9cm, inner sep=3pt, align=center},
  exlife/.style={dashed, black!45},
  exnote/.style={draw=black!40, fill=yellow!14, rounded corners=2pt, font=\scriptsize,
                 inner sep=2.5pt, align=center},
  exnum/.style={circle, text=white, font=\scriptsize\bfseries, inner sep=0pt, minimum size=0.46cm}
}
% \exarr{colour}{x from}{x to}{y}{frame number}{label}{sub-label}
\newcommand{\exarr}[7]{%
  \draw[-{Stealth[length=2.4mm]}, thick, #1] (#2,#4) -- (#3,#4);
  \node[exnum, fill=#1] at ($(#2,#4)!0.12!(#3,#4)$) {#5};
  \node[above=1pt, font=\footnotesize, text=#1, inner sep=1pt] at ($(#2,#4)!0.56!(#3,#4)$) {#6};
  \node[below=3pt, font=\scriptsize, text=black!75, inner sep=1pt, align=center]
       at ($(#2,#4)!0.5!(#3,#4)$) {#7};}

\subsection{Setup}
Table~\ref{tab:e-setup} lists the machines. The IP and MAC values come from \code{scenarios/proxy\_arp.json}. In this section, \code{.10} is short for 141.23.56.10, and \code{..:03:01} is short for \code{02:00:00:00:03:01}.

\begin{table}[htbp]
\centering
\caption{Machines in \code{proxy\_arp}. Every ARP cache starts empty, and so does the switch's MAC table.}
\label{tab:e-setup}
\footnotesize
\begin{tabularx}{\linewidth}{@{}l l l l l Y@{}}
\toprule
Node & LAN & IP / prefix & MAC & Gateway & Note \\
\midrule
h1 & A (bus) & 141.23.56.10/24 & \code{02:00:00:00:01:00} & .1 & sends the ping \\
h2 & A (bus) & 141.23.56.11/24 & \code{02:00:00:00:02:00} & .1 & bystander on LAN A \\
r1 if0 & A (bus) & 141.23.56.1/24 & \code{02:00:00:00:03:00} & --- & proxy list .21, .22, .23 \\
r1 if1 & B (star) & 141.23.56.17/29 & \code{02:00:00:00:03:01} & --- & router side of LAN B \\
h3 & B (star) & 141.23.56.21/29 & \code{02:00:00:00:04:00} & .17 & ping target \\
h4 & B (star) & 141.23.56.22/29 & \code{02:00:00:00:05:00} & .17 & bystander on LAN B \\
\bottomrule
\end{tabularx}
\end{table}

The ARP timers use the project defaults: an entry lives 30 s, and a pending request is repeated every 1 s. All times in this report are simulated time. The trace below uses a line delay of 100~ms, so each reply comes back before the first retry.

\subsection{Sequence diagram}
Figures~\ref{fig:e-seq} and~\ref{fig:e-seq2} show the 8 frames. Time runs downward. Each numbered arrow is one frame. A yellow note shows what a machine does with it, including cache changes. Machines that only listen are shown with grey notes. The colours are: orange for an ARP request, teal for an ARP reply, purple for the proxy reply, blue for IPv4 and ICMP.

\begin{figure}[htbp]
\centering
\begin{tikzpicture}[yscale=0.80]
\node[exhead, text=black!60, draw=black!50] (h2) at (0,0) {\textbf{h2}\\bystander};
\node[exhead] (h1) at (6,0) {\textbf{h1}\\.10};
\node[exhead, fill=routerfill] (r1) at (12,0) {\textbf{r1}\\proxy router};
\draw[exlife] (h2) -- (0,-10.1);
\draw[exlife] (h1) -- (6,-10.1);
\draw[exlife] (r1) -- (12,-10.1);
% step 1
\node[exnote, text width=4.4cm] at (6,-1.15) {ping .21: next hop is .21 itself,\\no cache entry: \code{pending},\\\code{queued}, \code{ip send}};
\exarr{arpreq}{6}{12}{-2.4}{1}{ARP request}{broadcast, tpa .21}
\node[exnote, text width=3.2cm, text=black!65, fill=black!5] at (0,-2.4) {hears frame 1: not for it,\\no cache entry, no learn};
% step 2
\node[exnote, text width=4.3cm] at (12,-3.85) {\code{learn} (if0,.10)=..:01:00\\\code{proxy\_reply}: .21 is in the list\\and leaves by if1};
\exarr{proxyc}{12}{6}{-5.3}{2}{ARP reply (proxy)}{unicast to ..:01:00, sha ..:03:00}
\node[exnote, text width=3.2cm, text=black!65, fill=black!5] at (0,-5.3) {frame 2: \code{rx ignored},\\not its MAC};
% step 3
\node[exnote, text width=4.4cm] at (6,-6.7) {\code{resolved} (if0,.21)=..:03:00\\the router's MAC; queue flushed};
\exarr{ipc}{6}{12}{-8.0}{3}{echo request}{TTL 64, to ..:03:00}
\node[exnote, text width=3.2cm, text=black!65, fill=black!5] at (0,-8.0) {frame 3: \code{rx ignored}};
\node[exnote, text width=4.3cm] at (12,-9.35) {TTL 64$\to$63, route via if1:\\\code{pending} (if1,.21), \code{queued},\\\code{ip forward} (continues below)};
\end{tikzpicture}
\caption{Frames 1 to 3, all on LAN A. The router answers the ARP request itself and h1 learns the router's MAC for .21.}
\label{fig:e-seq}
\end{figure}

\begin{figure}[htbp]
\centering
\begin{tikzpicture}[yscale=0.80]
\node[exhead] (h1) at (0,0) {\textbf{h1}\\.10};
\node[exhead, fill=routerfill] (r1) at (4.4,0) {\textbf{r1}\\if1 .17};
\node[exhead] (h3) at (8.8,0) {\textbf{h3}\\.21};
\node[exhead, text=black!60, draw=black!50] (h4) at (13.2,0) {\textbf{h4}\\bystander};
\foreach \x in {0,4.4,8.8,13.2} {\draw[exlife] (\x,-0.55) -- (\x,-14.6);}
% step 4
\exarr{arpreq}{4.4}{8.8}{-1.5}{4}{ARP request}{broadcast, tpa .21;\\switch learns ..:03:01, floods}
\node[exnote, text width=3.7cm, text=black!65, fill=black!5] at (13.2,-1.5) {hears frame 4: not for it,\\no learn};
% step 5
\node[exnote, text width=3.7cm] at (8.8,-3.1) {\code{learn} (if0,.17)=..:03:01\\\code{reply\_sent} for .21};
\exarr{arprep}{8.8}{4.4}{-4.6}{5}{ARP reply}{unicast to ..:03:01;\\switch learns ..:04:00}
\node[exnote, text width=3.7cm] at (4.4,-6.1) {\code{resolved} (if1,.21)=..:04:00\\queue flushed};
% step 6
\exarr{ipc}{4.4}{8.8}{-7.5}{6}{echo request}{TTL 63, unicast}
\node[exnote, text width=3.7cm] at (8.8,-9.0) {\code{ip deliver}, then \code{ip send}\\echo reply; gateway .17 is\\cached, so no ARP};
% step 7
\exarr{ipc}{8.8}{4.4}{-10.5}{7}{echo reply}{TTL 64, to ..:03:01}
\node[exnote, text width=3.7cm] at (4.4,-11.9) {\code{ip forward}, TTL 63\\(if0,.10) cached, no ARP};
% step 8
\exarr{ipc}{4.4}{0}{-13.3}{8}{echo reply}{TTL 63, to ..:01:00}
\node[exnote, text width=3.7cm] at (0,-14.6) {\code{ip deliver} echo\_rep:\\the ping is done};
\end{tikzpicture}
\caption{Frames 4 to 8. LAN B is a star, so unicast frames go only to their target. The router asks h3 for its MAC and then carries the ping and its reply.}
\label{fig:e-seq2}
\end{figure}

\subsection{Cache state over time}
Table~\ref{tab:e-cache} shows each cache just after each frame is sent. A cache key is (interface, IP address). Hosts h2 and h4 never learn anything, so their caches stay empty. The switch has its own MAC table. It learns ..:03:01 when frame 4 passes and ..:04:00 when frame 5 passes.

\begin{table}[htbp]
\centering
\caption{ARP cache contents just after each frame is sent. R means RESOLVED and P means PENDING.}
\label{tab:e-cache}
\footnotesize
\begin{tabularx}{\linewidth}{@{}l Y Y Y@{}}
\toprule
After frame & h1 & r1 & h3 \\
\midrule
start & empty & empty & empty \\
1 & (if0,.21) P & empty & empty \\
2 & (if0,.21) P & (if0,.10) = ..:01:00 R & empty \\
3 & (if0,.21) R = ..:03:00 & (if0,.10) = ..:01:00 R & empty \\
4 & (if0,.21) R = ..:03:00 & (if0,.10) R; (if1,.21) P & empty \\
5 & (if0,.21) R = ..:03:00 & (if0,.10) R; (if1,.21) P & (if0,.17) = ..:03:01 R \\
6 to 8 (end) & (if0,.21) R = ..:03:00 & (if0,.10) = ..:01:00 R; (if1,.21) = ..:04:00 R & (if0,.17) = ..:03:01 R \\
\bottomrule
\end{tabularx}
\end{table}

\subsection{What to notice}
\begin{enumerate}[leftmargin=1.6em, itemsep=2pt]
\item \textbf{The proxy reply carries the router's MAC.} In frame 2, the address pair is spa .21 with sha ..:03:00. So h1 maps .21 to the router, not to h3. h1 never learns h3's real MAC. This is the whole idea of proxy ARP (Pseudocode~\ref{alg:proxy}).
\item \textbf{The router lowers the TTL.} The echo request leaves h1 with TTL 64 and reaches h3 with 63. The echo reply starts at 64 and reaches h1 with 63. Each pass through r1 costs one.
\item \textbf{A second ping needs no ARP.} All three caches are full after the first ping. We ran the same preset with two pings, \code{runs/2026-11-16T18-14-26.jsonl}. The second echo request produced only the 4 ICMP frames and no ARP frame.
\item \textbf{A broadcast reaches h2 but does not teach it.} h2 gets frame 1, but it is neither the target nor a proxy, and it has no entry for h1. So it learns nothing. The same holds for h4 on LAN B.
\item \textbf{LAN A is a bus and LAN B is a star.} On the bus, h2 sees every frame and marks the unicast ones \code{ignored}. On the star, the switch has learned the MACs, so frames 5, 6 and 7 go only to r1 or h3, and h4 sees just the broadcast.
\end{enumerate}
